\section{Allocation as a Fixed Point}
\label{sec:allocation}

Suppose some packets are carried optically and never reach the radio. The pool
claim of \eqref{eq:demand} falls, because fewer vehicles contend; the
per-attempt failure probability $q$ falls with it; and $p_{\mathrm{RF}}$
improves for every packet that remains. But whether a packet is worth carrying
optically depends on $p_{\mathrm{RF}}$, which is what just changed. The
allocation and the channel it runs on determine each other, and neither can be
evaluated first.

This section states that circularity as a fixed point and solves it. The
resulting allocation is what Section~\ref{sec:results} evaluates against the
single-medium baselines.

\subsection{Coupling}

Let $f \in [0,1]$ denote the \emph{radio usage fraction}: the share of offered
packets that place a copy on the sidelink. A packet carried optically alone
returns its share of the pool, so three quantities fall together with $f$---the
number of contenders a transmitter competes with, the fraction of the period a
receiver spends transmitting rather than listening, and the pool claim. They are
one behaviour observed three ways and scale from one parameter. The per-attempt
failure probability is written $q(f)$ and the resulting per-packet radio risk
\begin{equation}
  p_{\mathrm{RF}}(f) = q(f)^{N} .
  \label{eq:rf_of_f}
\end{equation}

Half-duplex is why this coupling cannot be approximated away. At the attempt
counts of Table~\ref{tab:params} the half-duplex term is the same order as
collision, so charging a full duty cycle to a receiver that spent the period on
optics overstates radio failure by roughly $1/f$.

\subsection{Per-Packet Allocation}

Given the two risks, the allocation for a population of $K$ packets is the
cheapest assignment of media that meets the budget $\varepsilon$. Each packet
$k$ has a radio risk $p_{\mathrm{RF},k}$ and an optical risk
$p_{\mathrm{VLC},k}$, and may be assigned to the radio alone, the light alone,
or both. Activation costs are $1$, $1$ and $2$ respectively, so the cheapest
single-medium choice is whichever leg is less likely to fail,
\begin{equation}
  \pi_k = \min\!\left(p_{\mathrm{RF},k},\; p_{\mathrm{VLC},k}\right),
  \label{eq:single}
\end{equation}
and duplication is an upgrade that costs one further activation and, assuming
independence across media, reduces the packet's risk to
$p_{\mathrm{RF},k}\,p_{\mathrm{VLC},k}$.

Minimising mean cost subject to
$K^{-1}\sum_k \pi_k \le \varepsilon$ is then a linear relaxation whose solution
is greedy: baseline every packet on its cheaper leg, and upgrade packets to
duplication in decreasing order of the risk each upgrade removes,
\begin{equation}
  \Delta_k = \pi_k - p_{\mathrm{RF},k}\,p_{\mathrm{VLC},k},
  \label{eq:reduction}
\end{equation}
stopping as soon as the budget is met. Two properties of this rule matter
later. It can \emph{substitute}, not only duplicate---a packet whose optical
path is clean is routed to the light and off the radio entirely, which is what
relieves the pool---and it returns a mean activation cost, so the price of
meeting the budget is reported rather than assumed.

Independence across media is assumed in the $p_{\mathrm{RF}}p_{\mathrm{VLC}}$
term. The campaign measures the realised joint failure at a dependence ratio of
$1.18$ at $\rho = 30$, so the true joint outcome is slightly worse than
independence predicts and the duplication branch is correspondingly optimistic.
The allocation branch, which is the one that carries the result, does not use
that term.

\subsection{Solution}

Writing $\mathcal{A}$ for the allocation of \eqref{eq:single}--\eqref{eq:reduction}
and $f = \mathcal{A}\!\left(p_{\mathrm{RF}}(f), p_{\mathrm{VLC}}\right)$ for the
radio usage it induces, an operating point is a solution of
\begin{equation}
  f^{\star} = \mathcal{A}\!\left(p_{\mathrm{RF}}(f^{\star}),\, p_{\mathrm{VLC}}\right).
  \label{eq:fixedpoint}
\end{equation}
We solve \eqref{eq:fixedpoint} by damped iteration,
$f_{i+1} = (1-\lambda) f_i + \lambda \, \mathcal{A}(p_{\mathrm{RF}}(f_i), p_{\mathrm{VLC}})$,
with $\lambda = 0.25$.

Damping is not cosmetic. Undamped, the iteration oscillates at the higher
attempt counts: offloading improves the radio, which makes offloading
unattractive, which congests it again. That oscillation is an artefact of
updating every vehicle simultaneously---a real fleet decides asynchronously,
and $\lambda$ is what represents that.

Two properties of \eqref{eq:fixedpoint} require care.

\emph{The map has more than one fixed point.} Started from full radio use, a
congested pool makes duplication look necessary and duplicating keeps it
congested; the iteration settles there, self-consistently, at cost $2.0$ and a
miss rate that fails the budget. Started lower, the same inputs settle on most
packets going by light, the offered load collapsing, and the radio comfortably
carrying the remainder at cost $1.0$. Both are equilibria of the same system,
and at $\rho = 30$ on the pessimistic band they differ by a factor of $325$. We
therefore sweep several starting loads and report the cheapest feasible
equilibrium found, together with how many distinct equilibria were reached.
Bistability here is a property of the system rather than of the solver: a fleet
that begins congested stays congested absent something to coordinate the move
to light.

\emph{An exact fixed point need not exist.} The allocation is a discrete
ordering over a finite packet population, so $\mathcal{A}$ is a step function
and $\mathcal{A}(f) = f$ may have no solution. What exists is a small invariant
band. We therefore declare convergence on the reported quantities---the miss
rate and the mean cost---rather than on $f$ itself, which is what a published
number requires and what a stationarity test on a step function never delivers.

\subsection{Relation to Prior Fixed-Point Analyses}

Fixed-point analysis of C-V2X contention is established: Wang \emph{et al.}
solve one for semi-persistent scheduling, in which collision probability
depends on the idle resources a vehicle can sense, which in turn depend on the
collision probability~\cite{wang2018fixedpoint}. That fixed point is in the
collision probability \emph{at a fixed offered load}. The fixed point of
\eqref{eq:fixedpoint} is in the \emph{allocation}---in how much load leaves the
radio altogether---and it is the pool constraint \eqref{eq:demand}, absent from
single-medium analyses, that gives that departure its value.
